% Inactive proof fragment; uses the existing book macros and dependencies.
\subsection{Positive endpoints and right limits}

Fix a nonzero nonnegative $f\in L^1(\mathrm dV)$ on a connected compact Hermitian manifold $(X,h)$. Write $\Lambda_h(f)$ for the positive exponents admitting a genuine exact solution of
\[
 \mathrm{MA}^{\mathrm{ext}}_h(u)=\mathrm e^{\lambda u}f\,\mathrm dV.
\]
For $\lambda\in\Lambda_h(f)$, let $U_\lambda$ be the greatest solution supplied by \cref{thm:individual-exponential-subsolution}. The interval structure in \cref{prop:solvable-exponents-upper-interval} leaves the inclusion of a positive endpoint to be determined.

\begin{lemma}\label{lem:positive-exponent-uniform-upper-barrier}
For every $\tau>0$, there is a nonnegative $t_\tau\in\PSH(X,h)\cap C(X)$ such that every genuine exponential subsolution at an exponent $\lambda\geq\tau$ lies below $t_\tau$. Put $K_\tau=\sup_Xt_\tau$. For $\tau\leq\lambda_1\leq\lambda_2$ in $\Lambda_h(f)$,
\[
 U_{\lambda_2}\geq U_{\lambda_1}
             -K_\tau\frac{\lambda_2-\lambda_1}{\lambda_2}.
\]
Consequently $U_\lambda+K_\tau\log(\lambda/\tau)$ is increasing on $\Lambda_h(f)\cap[\tau,+\infty)$.
\end{lemma}
\begin{proof}
Set $q=\min(f,1)$. This is a bounded nonzero density, so \cref{thm:nguyen-exp} gives a continuous solution $w$ of $\mathrm{MA}_h(w)=\mathrm e^{\tau w}q\,\mathrm dV$. Put $t_\tau=w+\max(0,-\inf_Xw)$. Then $t_\tau\geq0$ and, for $\lambda\geq\tau$,
\[
 \mathrm{MA}_h(t_\tau)=\mathrm e^{\tau w}q\,\mathrm dV
       \leq\mathrm e^{\lambda t_\tau}f\,\mathrm dV.
\]
If $u$ is a genuine subsolution at $\lambda$, then on $\{t_\tau<u-\epsilon\}$,
\[
 \mathrm{MA}_h(t_\tau)
       \leq\mathrm e^{-\lambda\epsilon}
                  \mathrm{MA}^{\mathrm{ext}}_h(u-\epsilon).
\]
The bounded first potential admits a genuine $h/2$-envelope, and $u-\epsilon\preceq t_\tau$. Thus \cref{thm:singular-strict-domination} gives $t_\tau\geq u-\epsilon$. Let $\epsilon\downarrow0$.

For the order estimate, put $\delta=K_\tau(\lambda_2-\lambda_1)/\lambda_2$. Since $U_{\lambda_1}\leq K_\tau$, constant invariance gives
\[
 \mathrm{MA}^{\mathrm{ext}}_h(U_{\lambda_1}-\delta)
   =\mathrm e^{\lambda_1U_{\lambda_1}}f\,\mathrm dV
   \geq\mathrm e^{\lambda_2(U_{\lambda_1}-\delta)}f\,\mathrm dV.
\]
Greatest-subsolution maximality yields $U_{\lambda_2}\geq U_{\lambda_1}-\delta$. Finally $\log x\geq1-x^{-1}$ for $x\geq1$, which proves the logarithmically corrected order.
\end{proof}

\begin{proposition}\label{prop:positive-exponent-endpoint-attainment}
Suppose $\Lambda_h(f)\ne\varnothing$ and $a=\inf\Lambda_h(f)>0$. Then $a\in\Lambda_h(f)$ if and only if there are exponents $\lambda_j>a$ decreasing to $a$ and genuine exact solutions $u_j$ at $\lambda_j$ with $\inf_j\sup_Xu_j> -\infty$.

If $a\notin\Lambda_h(f)$, every sequence of genuine exact solutions at exponents $\lambda_j>a$ decreasing to $a$ satisfies
\[
 \sup_Xu_j\longrightarrow-\infty,\qquad
 \mathrm e^{\lambda_ju_j}\longrightarrow0\ \text{uniformly on }X,
 \qquad
 \|\mathrm{MA}^{\mathrm{ext}}_h(u_j)\|_{\mathrm{TV}}
                \longrightarrow0.
\]
The exponential has value zero at poles.
\end{proposition}
\begin{proof}
If $u_a$ is an exact solution at $a$, put $s=u_a-\max(0,\sup_Xu_a)$. Then $s\leq0$ and $\mathrm{MA}^{\mathrm{ext}}_h(s)\geq\mathrm e^{\lambda s}f\,\mathrm dV$ for every $\lambda\geq a$. Thus $U_\lambda\geq s$, which supplies the required sequence.

Conversely, the lemma with $\tau=a$ gives a common upper bound $K_a$ for $u_j$. Write $M_j=\sup_Xu_j$. Extract a subsequence for which $M_j\to M\in\mathbb R$ and, by \cref{thm:normalized-compactness}, $u_j-M_j\to v\in\PSH(X,h)$ in $L^1$. Then $u_j\to u=v+M$ in $L^1$, with a genuine limit and a common upper bound. After a further extraction the convergence holds almost everywhere. Since $\lambda_j\leq\lambda_1$ and $K_a\geq0$, dominated convergence gives
\[
 \mathrm e^{\lambda_ju_j}f\longrightarrow\mathrm e^{au}f
              \quad\text{in }L^1(\mathrm dV).
\]
Indeed all these densities are at most $\mathrm e^{\lambda_1K_a}f$. The lower-density closure \cref{lem:density-subsolution-lone-closure} gives $\mathrm{MA}^{\mathrm{ext}}_h(u)\geq\mathrm e^{au}f\,\mathrm dV$. The greatest-solution theorem now supplies an exact solution at $a$.

If $a$ is not admitted and the suprema of an exact solution sequence failed to tend to $-\infty$, a subsequence would satisfy the lower-bound condition just proved. Thus $M_j\to-\infty$. Eventually $M_j<0$; since $\lambda_j\geq a>0$,
\[
 0\leq\mathrm e^{\lambda_ju_j}\leq\mathrm e^{aM_j}
                 \longrightarrow0.
\]
The exact mass is at most $\mathrm e^{aM_j}\|f\|_1$, proving the last assertion.
\end{proof}

\begin{proposition}\label{prop:greatest-exponential-right-continuity}
For every $a\in\Lambda_h(f)$, the greatest solutions satisfy
\[
 U_\lambda\longrightarrow U_a
       \quad\text{in }L^1\text{ and in capacity}\quad(\lambda\downarrow a).
\]
Their Monge--Amp\`ere measures converge in total variation to $\mathrm e^{aU_a}f\,\mathrm dV$.
\end{proposition}
\begin{proof}
Put $K=K_a$ from the lemma. For any $\lambda_j>a$ decreasing to $a$, the potentials
\[
 V_j=U_{\lambda_j}+K\log(\lambda_j/a)
\]
decrease and lie above $U_a$. Their limit $V$ is genuine. Dominated convergence between the integrable bounds $U_a$ and $V_1$ gives $L^1$ convergence, and \cref{lem:monotone-psh-capacity} gives capacity convergence. Since the added constants tend to zero, $U_{\lambda_j}\to V$ in both topologies.

Apply \cref{thm:exponential-closedness} to the nonpositive potentials $U_{\lambda_j}-K$, the exponents $\lambda_j\to a>0$ and the densities $\mathrm e^{\lambda_jK}f\to\mathrm e^{aK}f$ in $L^1$. It gives $\mathrm{MA}^{\mathrm{ext}}_h(V)=\mathrm e^{aV}f\,\mathrm dV$, with total-variation convergence of the equation measures. Hence $V\leq U_a$ by greatest-solution maximality, while the lower bound gives $V\geq U_a$. Thus $V=U_a$. The argument applies to every decreasing sequence of exponents, giving the stated right limit.
\end{proof}

The endpoint test is expressed in terms of existing solutions. An intrinsic criterion on $f$, the existence or nonexistence of a positive excluded endpoint, and continuity from the left remain separate questions. The right limit selects greatest solutions; it does not identify limits of arbitrary solutions with other singularities.
